\section{Model Shift-Reduce, Stack, dan Handle}
\label{sec:ch11-shift-reduce}

Model shift-reduce pada bagian ini mengikuti Thain Bab 4.4.1 dan
Mogensen Bab 3.14.

\subsection{Ide Dasar}

Parser bottom-up membangun pohon urai dari daun ke akar. Pada setiap
langkah, parser membaca token dari kiri ke kanan dan memutuskan
apakah akan \eng{shift} (mendorong token ke stack) atau
\eng{reduce} (mengganti ruas kanan aturan dengan nonterminal).
Parser menerima masukan ketika stack berisi simbol awal dan seluruh
masukan telah dibaca.

\begin{figure}[H]
\centering
\begin{tikzpicture}[node distance=0mm]
  \node[memorycell, minimum width=3.0cm, minimum height=2.6cm,
        align=center, fill=ubbblue!5] (stack) {Stack\\\(\vdots\)\\handle};
  \node[memorycell, minimum width=4.0cm, minimum height=2.6cm,
        right=6mm of stack, align=center, fill=ubbteal!5] (input)
        {Masukan\\\(\code{id} + \code{id} * \code{id}\ \$\)};
  \node[phase, minimum width=2.8cm, right=6mm of input] (action)
        {Aksi:\\shift/reduce};
  \draw[arrow] (stack) -- (action);
  \draw[arrow] (input) -- (action);
\end{tikzpicture}
\caption{Model shift-reduce dengan stack dan masukan}
\label{fig:ch11-model}
\end{figure}

\subsection{Contoh Tata Bahasa}

Sepanjang bab ini digunakan tata bahasa ekspresi:
\[
E \to E + T \mid T, \quad
T \to T * F \mid F, \quad
F \to \code{id}.
\]
Tata bahasa ini left-recursive, sehingga cocok untuk bottom-up tetapi
tidak cocok untuk top-down tanpa transformasi.

\subsection{Handle}

\begin{definisi}[Handle]
Handle adalah substring pada bentuk kalimat yang cocok dengan ruas
kanan suatu aturan produksi dan dapat direduksi menjadi nonterminal
sebelah kiri.
\end{definisi}

Pada bentuk kalimat \(\code{id} + \code{id} * \code{id}\), handle
pertama adalah \code{id} yang direduksi menjadi \(F\), lalu \(T\),
lalu \(E\). Pemilihan handle yang tepat menentukan keberhasilan
parsing.

\subsection{Penelusuran Shift-Reduce}

Untuk masukan \(\code{id} + \code{id} * \code{id}\ \$\):
\begin{table}[H]
\centering
\caption{Penelusuran shift-reduce}
\label{tab:ch11-trace}
\begin{tabularx}{\textwidth}{@{}l X l@{}}
\toprule
\textbf{Stack} & \textbf{Masukan} & \textbf{Aksi} \\
\midrule
 & \(\code{id} + \code{id} * \code{id}\ \$\) & shift \\
\(\code{id}\) & \(+ \code{id} * \code{id}\ \$\) &
reduce \(F \to \code{id}\) \\
\(F\) & \(+ \code{id} * \code{id}\ \$\) & reduce \(T \to F\) \\
\(T\) & \(+ \code{id} * \code{id}\ \$\) & reduce \(E \to T\) \\
\(E\) & \(+ \code{id} * \code{id}\ \$\) & shift \(+\) \\
\(E+\) & \(\code{id} * \code{id}\ \$\) & shift \code{id} \\
\(E+\code{id}\) & \(* \code{id}\ \$\) & reduce \(F \to \code{id}\) \\
\(E+F\) & \(* \code{id}\ \$\) & reduce \(T \to F\) \\
\(E+T\) & \(* \code{id}\ \$\) & shift \(*\) \\
\(E+T*\) & \(\code{id}\ \$\) & shift \code{id} \\
\(E+T*\code{id}\) & \(\$\) & reduce \(F \to \code{id}\) \\
\(E+T*F\) & \(\$\) & reduce \(T \to T*F\) \\
\(E+T\) & \(\$\) & reduce \(E \to E+T\) \\
\(E\) & \(\$\) & terima \\
\bottomrule
\end{tabularx}
\end{table}

Penelusuran pada \tabref{tab:ch11-trace} menunjukkan bahwa parser
tidak perlu melihat seluruh masukan. Keputusan diambil berdasarkan
isi stack dan token saat ini.

\subsection{Keuntungan}

Bottom-up parsing dapat menangani left recursion secara alami,
mendukung grammar yang lebih luas, dan menjadi dasar parser LR.
Keterbatasannya adalah tabel dan algoritma yang lebih kompleks
daripada recursive descent.
